\item \model{GLM-5}

\paragraph*{گذار از توجه تنک موضعی به توجه شاخص‌گذاری‌شده سراسری}
در گزارش فنی مدل \model{GLM-5} \cite{glm2026glm5}، چالش مقیاس‌پذیری مکانیزم موقعیت در امتداد بسط پنجره بافت از ۴ هزار به ۲۰۰ هزار توکن (\en{4K to 200K}) در فاز آموزش میانی مورد واکاوی جامع قرار گرفته است. تحلیل رفتار تجربی نشان می‌دهد که جایگزینی کدگذاری موقعیت متراکم با روش‌های توجه برشی ثابت نظیر توجه پنجره لغزان همگن (\en{SWA Interleave}) منجر به فروپاشی فاجعه‌بار در درک موقعیتی بلندمدت می‌گردد؛ به‌طوری‌که در ارزیابی معیار \en{RULER@128K}، مدل با افت ۳۰.۳۵ نمره‌ای مواجه می‌شود (کاهش از ۷۵.۲۸ به ۴۴.۹۳).

\paragraph*{حفظ تمامیت میدان توجه بدون افت دقت فاز}
برای برطرف ساختن ناپایداری‌های فازی و بازیابی وابستگی‌های دوربرد در فضای RoPE، \model{GLM-5} رویکرد توجه تنک دیپ‌سیک (\en{DSA}) را اتخاذ کرده است. شاخص‌گذار صاعقه‌ای (\en{Lightning Indexer}) به صورت پویا و محتوا-محور، توکن‌های پراهمیت در سرتاسر بافت ۲۰۰K را بدون حذف فواصل بلندمدت انتخاب می‌کند:
\begin{equation}
    I_{t, s} = \sum_{h=1}^{n_h^I} w_{t, h}^I \cdot \text{ReLU}\left( \langle q_{t, h}^I, k_s^I \rangle \right)
\end{equation}
از منظر آماری و اطلاعاتی، چون ساختار گزینش توکن در DSA از پیش‌فرض‌های هندسی متقارن پنجره‌ای مستقل است، پیوستگی فضایی بردارهای چرخانده‌شده توسط RoPE در کل لایه‌ها بدون نیاز به بازنشانی مقیاس فرکانس حفظ می‌گردد. همچنین در طول مرحله آموزش میانی، جهت پایدارسازی درک فواصل دور و مهار پدیده «گم‌شدن در میانه» (\en{Lost-in-the-Middle})، داده‌های مصنوعی با وابستگی‌های ترتیبی بلند (\en{NextLong} و \en{EntropyLong}) تزریق شده‌اند تا گرادیان‌های فازی در فواصل انتهایی پنجره ۲۰۰K از محوشدگی مصون بمانند.